\FloatBarrier
\section{Downstream joint clearance revision}
\label{sec:downstream-joints}

A subsequent revision resolves the remaining contacts between Downstream\_Beam\_tube\_Segment\_1 and both
Beamline\_Adapter\_1 and Downstream\_Beam\_tube\_Segment\_2. The nominal cylindrical gaps are 0.0508 and
0.0127 mm, but the mating surfaces use 24-sided polygons with different angular alignments. Both alignments
are present in the engineering inputs. The reported 0.0398313 and 0.0568888 mm Geant4 penetration samples
are not global maxima. Polygon support planes give deficits of 0.342860 and 0.395184 mm at the respective
joints. Independent triangle checks find 240 and 239 contacts in the previous meshes; positive
vertex-distance samples alone would not establish clearance.

Only Segment 1 is changed. Its entrance socket radius increases from 60.325 to 60.771702 mm through nominal
z = 5048.376367 mm, followed by a 1 mm transition back to the original bore. Two additional 24-vertex rings
replace one long wall band with three bands, adding 96 triangles. The downstream socket radius increases
from 70.6628 to 71.162277 mm over its existing downstream wall interval. Both sockets retain their angular
grids. The adapter, Segment 2, all outer bounds, materials, passive placements, and the open beam passage
remain unchanged.

The revised Segment 1 contains 5,096 triangles and 254,884 bytes, compared with 5,000 triangles and 250,084
bytes before relief. It remains one closed manifold genus-one component with no detected self-intersections
or zero-area triangles. Final volume is 6710837.897365 mm$^3$, removing 0.604764\% steel relative to the
previous closed mesh. Volume change against the original annularly repaired reference is -0.620987\%. The
original engineering surface remains open and has no valid mass. Maximum sampled distance from that source
remains 0.745438 mm, within the existing 0.75 mm guard. The current builder reproduces the relieved STL byte
for byte.

Independent clipping of neighboring triangles to their shared axial intervals gives minimum clearance to all
48 socket-wall triangle planes of 0.100019 mm at the adapter and 0.100022 mm at Segment 2. Triangle
collisions and bidirectional vertex containment are zero for both pairs. All eleven source CAD
configurations containing the three volumes use the same identity placements, so the same joint checks apply
to each. The separate 10,000-point Geant4 assembly test gives eight overlap reports before relief and no
exceptions afterward. Full default beamline construction now finishes without missing meshes, solid
warnings, or overlap reports. These checks cover the tested configuration and do not certify every possible
complete detector assembly.

\begin{table}[htbp]
\centering
\caption{Segment 1 before and after socket relief.}
\label{tab:downstream-joints-mesh}
\setlength{\tabcolsep}{9pt}
\begin{tabular}{lrrr}
\toprule
Quantity & Previous & Relieved & Change \\
\midrule
Triangles & 5,000 & 5,096 & +96 \\
Bytes & 250,084 & 254,884 & +4,800 \\
Boundary edges & 0 & 0 & 0 \\
Volume (mm$^3$) & 6751669.593 & 6710837.897 & -0.604764\% \\
\bottomrule
\end{tabular}
\end{table}

A new tracking comparison includes the adapter and both downstream segments in air. Its baseline is the
previous optimized closed meshes, not the original engineering inputs. The fixed GEMC3 executable,
materials, identity placements, matched 5 GeV target-origin primaries over 0.1--1.3 degrees and full
azimuth, FTFP\_BERT, no optics or field, one thread, default production cuts, and disabled GUI/event output
are identical. Five alternating-order repetitions use seeds 12345--12349 and 200 events per electron,
proton, positive-pion, and neutron run. All 52 processes finish normally and all forty transport runs report
200 events.

Previous/relieved median transport ratios are 0.981, 1.000, 1.008, 0.960 for electrons, protons, positive
pions, and neutrons respectively. Ratios above one denote faster tracking. The previous assembly produces 18
GeomNav1002, 4 GeomVol1002 exception blocks across the recorded processes; the relieved assembly produces 15
GeomNav1002 blocks. The remaining navigation warnings occur in the unchanged bellows region rather than
either relieved joint. Their z coordinates lie between 10177.906421 and 10281.864247 mm; all 1,741 triangles
intersecting the combined before/after warning interval are identical. Clean overlap checks therefore do not
establish warning-free tracking. These timings include overlap removal, local steel removal, and 96
additional facets. They measure the repair workload rather than a pure decimation gain, and do not establish
full CLAS12 performance or physics-response equivalence. The earlier study chapters and archives retain
their historical measurements.

\begin{table}[htbp]
\centering
\caption{Median transport wall seconds per 200 events: three-part downstream assembly.}
\label{tab:downstream-joints-tracking}
\setlength{\tabcolsep}{9pt}
\begin{tabular}{lrrr}
\toprule
Particle & Previous & Relieved & Previous/relieved \\
\midrule
e- & 11.25010 & 11.46960 & 0.981 \\
proton & 1.75729 & 1.75794 & 1.000 \\
pi+ & 1.93580 & 1.92101 & 1.008 \\
neutron & 1.09550 & 1.14060 & 0.960 \\
\bottomrule
\end{tabular}
\end{table}

\FloatBarrier
\begin{figure}[htbp]
\centering
\includegraphics[width=0.95\textwidth]{figures/downstream_joint_clearance.pdf}
\caption{Actual socket and neighboring tube polygons projected onto the most restrictive socket face at each
joint. The previous socket cuts into its neighbor; the relieved socket provides at least 0.1 mm face
clearance while the neighboring tube remains unchanged.}
\label{fig:downstream-joints}
\end{figure}
